
\documentclass{article}
\usepackage{colortbl}
\usepackage{makecell}
\usepackage{multirow}
\usepackage{supertabular}

\begin{document}

\newcounter{utterance}

\centering \large Interaction Transcript for game `dond', experiment `coop\_en', episode 14 with refineonly{-}Qwen3.5{-}27B.
\vspace{24pt}

{ \footnotesize  \setcounter{utterance}{1}
\setlength{\tabcolsep}{0pt}
\begin{supertabular}{c@{$\;$}|p{.15\linewidth}@{}p{.15\linewidth}p{.15\linewidth}p{.15\linewidth}p{.15\linewidth}p{.15\linewidth}}
    \# & $\;$A & \multicolumn{4}{c}{Game Master} & $\;\:$B\\
    \hline

    \theutterance \stepcounter{utterance}  
    & & \multicolumn{4}{p{0.6\linewidth}}{
        \cellcolor[rgb]{0.9,0.9,0.9}{
            \makecell[{{p{\linewidth}}}]{
                \texttt{\tiny{[P1$\langle$GM]}}
                \texttt{You are playing a cooperative negotiation game in which you have to agree on how to divide a set of items among you and another player.} \\
\\ 
\texttt{Rules:} \\
\texttt{(a) You and the other player are given a set of items. Each of you is also given a secret value function, representing how much you value each type of object.} \\
\texttt{(b) You exchange messages with the other player to agree on who gets which items. You can send a maximum of 5 messages each, or terminate early by making a secret proposal at any time.} \\
\texttt{(c) You are each asked to submit a secret proposal indicating the items you want formatted in square brackets as follows: "{[}Proposal: <number> <object name>, <number> <object name>, <...>{]}"} \\
\texttt{(d) If your proposals are complementary, i.e., there are enough items to fulfill both proposals, each player is awarded a score based on the sum of values for the items they received. Otherwise, both of you get zero points.} \\
\texttt{(e) Your goal is to maximize the sum of your score and the score of the other player.} \\
\\ 
\texttt{Let us start.} \\
\\ 
\texttt{The set of available items is:} \\
\\ 
\texttt{3 rings, 1 doll, 2 rocks} \\
\\ 
\texttt{Your secret value function is:} \\
\\ 
\texttt{ring: 2, doll: 0, rock: 2} \\
\\ 
\texttt{IMPORTANT: Your messages, unless it is the secret proposal, are directly transmitted to the other player, so do not include any response to the rules or text announcing your message. To make a secret proposal, use the indicated format. Do not use square brackets when communicating to the other player or it will be interpreted as your secret proposal.} \\
            }
        }
    }
    & \\ \\

    \theutterance \stepcounter{utterance}  
    & \multicolumn{4}{p{0.6\linewidth}}{
        \cellcolor[rgb]{0.9,0.9,0.9}{
            \makecell[{{p{\linewidth}}}]{
                \texttt{\tiny{[P1$\rangle$GM]}}
                \texttt{I value rings and rocks equally at 2 points each, and I have no use for the doll (0 points). Since there are 3 rings and 2 rocks available, and I need to maximize our combined score, I suggest we split the items based on our preferences. Could you share your value function for the items? That way, we can figure out who should take the doll and how to best divide the rings and rocks to ensure we both get high scores.} \\
            }
        }
    }
    & & \\ \\

    \theutterance \stepcounter{utterance}  
    & & & \multicolumn{2}{p{0.3\linewidth}}{
        \cellcolor[rgb]{0.9,0.9,0.9}{
            \makecell[{{p{\linewidth}}}]{
                \texttt{\tiny{[GM$|$GM]}}
                \texttt{continue} \\
            }
        }
    }
    & & \\ \\

    \theutterance \stepcounter{utterance}  
    & & & \multicolumn{2}{p{0.3\linewidth}}{
        \cellcolor[rgb]{0.9,0.9,0.9}{
            \makecell[{{p{\linewidth}}}]{
                \texttt{\tiny{[GM$|$GM]}}
                \texttt{I value rings and rocks equally at 2 points each, and I have no use for the doll (0 points). Since there are 3 rings and 2 rocks available, and I need to maximize our combined score, I suggest we split the items based on our preferences. Could you share your value function for the items? That way, we can figure out who should take the doll and how to best divide the rings and rocks to ensure we both get high scores.} \\
            }
        }
    }
    & & \\ \\

    \theutterance \stepcounter{utterance}  
    & & \multicolumn{4}{p{0.6\linewidth}}{
        \cellcolor[rgb]{0.9,0.9,0.9}{
            \makecell[{{p{\linewidth}}}]{
                \texttt{\tiny{[GM$\rangle$P2]}}
                \texttt{You are playing a cooperative negotiation game in which you have to agree on how to divide a set of items among you and another player.} \\
\\ 
\texttt{Rules:} \\
\texttt{(a) You and the other player are given a set of items. Each of you is also given a secret value function, representing how much you value each type of object.} \\
\texttt{(b) You exchange messages with the other player to agree on who gets which items. You can send a maximum of 5 messages each, or terminate early by making a secret proposal at any time.} \\
\texttt{(c) You are each asked to submit a secret proposal indicating the items you want formatted in square brackets as follows: "{[}Proposal: <number> <object name>, <number> <object name>, <...>{]}"} \\
\texttt{(d) If your proposals are complementary, i.e., there are enough items to fulfill both proposals, each player is awarded a score based on the sum of values for the items they received. Otherwise, both of you get zero points.} \\
\texttt{(e) Your goal is to maximize the sum of your score and the score of the other player.} \\
\\ 
\texttt{Let us start.} \\
\\ 
\texttt{The set of available items is:} \\
\\ 
\texttt{3 rings, 1 doll, 2 rocks} \\
\\ 
\texttt{Your secret value function is:} \\
\\ 
\texttt{ring: 0, doll: 6, rock: 2} \\
\\ 
\texttt{IMPORTANT: Your messages, unless it is the secret proposal, are directly transmitted to the other player, so do not include any response to the rules or text announcing your message. To make a secret proposal, use the indicated format. Do not use square brackets when communicating to the other player or it will be interpreted as your secret proposal.} \\
\\ 
\\ 
\texttt{I value rings and rocks equally at 2 points each, and I have no use for the doll (0 points). Since there are 3 rings and 2 rocks available, and I need to maximize our combined score, I suggest we split the items based on our preferences. Could you share your value function for the items? That way, we can figure out who should take the doll and how to best divide the rings and rocks to ensure we both get high scores.} \\
            }
        }
    }
    & \\ \\

    \theutterance \stepcounter{utterance}  
    & & & \multicolumn{4}{p{0.6\linewidth}}{
        \cellcolor[rgb]{0.9,0.9,0.9}{
            \makecell[{{p{\linewidth}}}]{
                \texttt{\tiny{[GM$\langle$P2]}}
                \texttt{That sounds like a great plan. Since you value rings and rocks at 2 points each and have no use for the doll, and I value the doll at 6 points (with rings at 0 and rocks at 2), the most efficient way to maximize our combined score is for me to take the doll.} \\
\\ 
\texttt{For the remaining items, since you value rings and rocks equally, and I value rocks at 2 (same as you) but rings at 0, it makes sense for you to take all the rings and rocks. Since I get 0 from rings, you should definitely take all 3 rings. For the rocks, since we both value them at 2, it doesn't matter who gets them for the total score, but to keep it simple, you could take the rocks as well.} \\
\\ 
\texttt{So, a proposed split would be:} \\
\texttt{{-} You get: 3 rings, 2 rocks} \\
\texttt{{-} I get: 1 doll} \\
\\ 
\texttt{This way, you get 10 points, and I get 6 points, totaling 16 points, which is the maximum possible.} \\
\\ 
\texttt{Does this work for you? If so, I can make a secret proposal.} \\
            }
        }
    }
     \\ \\

    \theutterance \stepcounter{utterance}  
    & & & \multicolumn{2}{p{0.3\linewidth}}{
        \cellcolor[rgb]{0.9,0.9,0.9}{
            \makecell[{{p{\linewidth}}}]{
                \texttt{\tiny{[GM$|$GM]}}
                \texttt{continue} \\
            }
        }
    }
    & & \\ \\

    \theutterance \stepcounter{utterance}  
    & & & \multicolumn{2}{p{0.3\linewidth}}{
        \cellcolor[rgb]{0.9,0.9,0.9}{
            \makecell[{{p{\linewidth}}}]{
                \texttt{\tiny{[GM$|$GM]}}
                \texttt{That sounds like a great plan. Since you value rings and rocks at 2 points each and have no use for the doll, and I value the doll at 6 points (with rings at 0 and rocks at 2), the most efficient way to maximize our combined score is for me to take the doll.} \\
\\ 
\texttt{For the remaining items, since you value rings and rocks equally, and I value rocks at 2 (same as you) but rings at 0, it makes sense for you to take all the rings and rocks. Since I get 0 from rings, you should definitely take all 3 rings. For the rocks, since we both value them at 2, it doesn't matter who gets them for the total score, but to keep it simple, you could take the rocks as well.} \\
\\ 
\texttt{So, a proposed split would be:} \\
\texttt{{-} You get: 3 rings, 2 rocks} \\
\texttt{{-} I get: 1 doll} \\
\\ 
\texttt{This way, you get 10 points, and I get 6 points, totaling 16 points, which is the maximum possible.} \\
\\ 
\texttt{Does this work for you? If so, I can make a secret proposal.} \\
            }
        }
    }
    & & \\ \\

    \theutterance \stepcounter{utterance}  
    & & \multicolumn{4}{p{0.6\linewidth}}{
        \cellcolor[rgb]{0.9,0.9,0.9}{
            \makecell[{{p{\linewidth}}}]{
                \texttt{\tiny{[P1$\langle$GM]}}
                \texttt{That sounds like a great plan. Since you value rings and rocks at 2 points each and have no use for the doll, and I value the doll at 6 points (with rings at 0 and rocks at 2), the most efficient way to maximize our combined score is for me to take the doll.} \\
\\ 
\texttt{For the remaining items, since you value rings and rocks equally, and I value rocks at 2 (same as you) but rings at 0, it makes sense for you to take all the rings and rocks. Since I get 0 from rings, you should definitely take all 3 rings. For the rocks, since we both value them at 2, it doesn't matter who gets them for the total score, but to keep it simple, you could take the rocks as well.} \\
\\ 
\texttt{So, a proposed split would be:} \\
\texttt{{-} You get: 3 rings, 2 rocks} \\
\texttt{{-} I get: 1 doll} \\
\\ 
\texttt{This way, you get 10 points, and I get 6 points, totaling 16 points, which is the maximum possible.} \\
\\ 
\texttt{Does this work for you? If so, I can make a secret proposal.} \\
            }
        }
    }
    & \\ \\

    \theutterance \stepcounter{utterance}  
    & \multicolumn{4}{p{0.6\linewidth}}{
        \cellcolor[rgb]{0.9,0.9,0.9}{
            \makecell[{{p{\linewidth}}}]{
                \texttt{\tiny{[P1$\rangle$GM]}}
                \texttt{That sounds perfect. Your analysis is correct: you take the doll for 6 points, and I take the 3 rings and 2 rocks for 10 points, maximizing our total score to 16.} \\
\\ 
\texttt{{[}Proposal: 3 rings, 2 rocks{]}} \\
            }
        }
    }
    & & \\ \\

    \theutterance \stepcounter{utterance}  
    & & & \multicolumn{2}{p{0.3\linewidth}}{
        \cellcolor[rgb]{0.9,0.9,0.9}{
            \makecell[{{p{\linewidth}}}]{
                \texttt{\tiny{[GM$|$GM]}}
                \texttt{proposal} \\
            }
        }
    }
    & & \\ \\

    \theutterance \stepcounter{utterance}  
    & & & \multicolumn{2}{p{0.3\linewidth}}{
        \cellcolor[rgb]{0.9,0.9,0.9}{
            \makecell[{{p{\linewidth}}}]{
                \texttt{\tiny{[GM$|$GM]}}
                [3, 0, 2]
            }
        }
    }
    & & \\ \\

    \theutterance \stepcounter{utterance}  
    & & \multicolumn{4}{p{0.6\linewidth}}{
        \cellcolor[rgb]{0.9,0.9,0.9}{
            \makecell[{{p{\linewidth}}}]{
                \texttt{\tiny{[GM$\rangle$P2]}}
                \texttt{The other player has submitted a secret proposal. It is now your turn to submit a proposal of your own.} \\
            }
        }
    }
    & \\ \\

    \theutterance \stepcounter{utterance}  
    & & & \multicolumn{4}{p{0.6\linewidth}}{
        \cellcolor[rgb]{0.9,0.9,0.9}{
            \makecell[{{p{\linewidth}}}]{
                \texttt{\tiny{[GM$\langle$P2]}}
                \texttt{{[}Proposal: 3 rings, 2 rocks{]}} \\
            }
        }
    }
     \\ \\

    \theutterance \stepcounter{utterance}  
    & & & \multicolumn{2}{p{0.3\linewidth}}{
        \cellcolor[rgb]{0.9,0.9,0.9}{
            \makecell[{{p{\linewidth}}}]{
                \texttt{\tiny{[GM$|$GM]}}
                \texttt{proposal} \\
            }
        }
    }
    & & \\ \\

    \theutterance \stepcounter{utterance}  
    & & & \multicolumn{2}{p{0.3\linewidth}}{
        \cellcolor[rgb]{0.9,0.9,0.9}{
            \makecell[{{p{\linewidth}}}]{
                \texttt{\tiny{[GM$|$GM]}}
                [3, 0, 2]
            }
        }
    }
    & & \\ \\

    \theutterance \stepcounter{utterance}  
    & & & \multicolumn{2}{p{0.3\linewidth}}{
        \cellcolor[rgb]{0.9,0.9,0.9}{
            \makecell[{{p{\linewidth}}}]{
                \texttt{\tiny{[GM$|$GM]}}
                [[3, 0, 2], [3, 0, 2]]
            }
        }
    }
    & & \\ \\

\end{supertabular}
}

\end{document}
